\section{Einführung: Warum die Masse des Photons von Bedeutung ist}

Die Frage, ob das Photon tatsächlich masselos ist, ist keineswegs nur 
philosophischer Natur. Sie berührt die Grundlagen der Elektrodynamik, 
die interne Konsistenz quantenfeldtheoretischer Beschreibungen und die 
Grenzen unserer experimentellen Möglichkeiten. Die klassische 
Elektrodynamik setzt ein masseloses Photon voraus, doch diese Annahme 
sollte nicht unkritisch übernommen werden: Bereits eine extrem kleine, 
aber von Null verschiedene Photonenmasse hätte klar nachweisbare 
physikalische Konsequenzen.

Eine endliche Photonenmasse würde die Maxwell-Gleichungen verändern, die 
Struktur elektromagnetischer Felder modifizieren und das langreichweitige 
Verhalten der Coulomb-Wechselwirkung aufheben. Die entsprechende Proca-Theorie 
sagt eine endliche Reichweite des elektromagnetischen Feldes, einen 
longitudinalen Polarisationsmodus und veränderte Ausbreitungseigenschaften 
elektromagnetischer Strahlung voraus. Nichts davon wurde experimentell 
beobachtet, doch die theoretischen Konsequenzen erlauben es, präzise 
Grenzen für eine mögliche Photonenmasse abzuleiten.

Die moderne Physik bietet mehrere unabhängige Zugänge zu dieser Frage. 
Astrophysikalische Messungen untersuchen elektromagnetische Felder über 
galaktische und interplanetare Distanzen hinweg und liefern Obergrenzen, 
die weit über die Möglichkeiten von Laborexperimenten hinausgehen. 
Innerhalb der Quantenelektrodynamik ist die Masselosigkeit des Photons 
eng mit der $U(1)$-Eichsymmetrie und der Renormierbarkeit der Theorie 
verknüpft. Selbst das Verhalten verschränkter Photonen offenbart die 
Struktur eines masselosen Spin-1-Teilchens: Das experimentelle 
Zweizustandsverhalten ist nur mit einem masselosen Eichboson vereinbar.

Dieses Kapitel skizziert die Motivation, die Photonenmasse aus einer 
vereinheitlichten Perspektive zu betrachten. Die folgenden Abschnitte 
führen Symmetrieprinzipien, feldtheoretische Argumente, astrophysikalische 
Daten, Präzisionsmessungen sowie die Helizitätsstruktur quantenmechanischer 
Verschränkung zu einem kohärenten Gesamtbild zusammen: Falls das Photon 
eine Masse besitzt, muss sie außerordentlich klein sein. Ziel ist es nicht 
nur, die aktuellen Grenzen darzustellen, sondern zu verdeutlichen, warum 
die Masselosigkeit eine natürliche und stabile Eigenschaft im theoretischen 
Rahmen der Elektrodynamik ist.

